% Generated draft: requires mathematical and source review.
\documentclass[11pt]{article}
\usepackage[margin=25mm]{geometry}
\usepackage{fontspec}
\setmainfont{Latin Modern Roman}
\usepackage{amsmath,amssymb,mathtools,unicode-math}
\setmathfont{Latin Modern Math}
\usepackage{enumitem,needspace,xurl,hyperref}
\setlength{\parindent}{0pt}
\setlength{\parskip}{5pt}
\setlength{\emergencystretch}{3em}
\newcommand{\N}{\mathbb{N}}
\newcommand{\Z}{\mathbb{Z}}
\newcommand{\Q}{\mathbb{Q}}
\newcommand{\R}{\mathbb{R}}
\newcommand{\C}{\mathbb{C}}
\newcommand{\F}{\mathbb{F}}
\newcommand{\A}{\mathbb{A}}
\newcommand{\PP}{\mathbb{P}}
\newcommand{\E}{\mathbb{E}}
\newcommand{\eps}{\varepsilon}
\newcommand{\dd}{\,\mathrm{d}}
\begin{document}
\section*{UnsolvedMath ID: 801045}
\textbf{Source problem code:} catalogue entry
\subsection{Definitions and mathematical statement}
Let \(X\) be a standard Borel space and let \(\vec G\subseteq X\times X\) be a Borel directed graph. Local finiteness means that, for every \(x\in X\), the set \(\{y\in X:(x,y)\in\vec G\text{ or }(y,x)\in\vec G\}\) is finite. Uniformly bounded in-degrees means that there exists an integer \(d\geq 0\) such that
\[
\bigl|\{y\in X:(y,x)\in\vec G\}\bigr|\leq d\qquad(x\in X).
\]
For an integer \(n\geq 0\), define the forward ball
\[
B^+_{\vec G}(x,n)=\{y\in X:\text{there is a directed path from }x\text{ to }y\text{ of length at most }n\},
\]
including the length-zero path. Define the connected equivalence relation \(E_{\vec G}\) by declaring \(x\,E_{\vec G}\,y\) when there is a finite path from \(x\) to \(y\) after edge directions are forgotten. A Borel equivalence relation \(E\) is hyperfinite if there are Borel equivalence relations \(E_m\), indexed by integers \(m\geq 0\), each with finite equivalence classes, such that
\[
E_m\subseteq E_{m+1}\quad\text{and}\quad E=\bigcup_{m\geq 0}E_m.
\]
The supplied question asks: if \(\vec G\) is Borel and locally finite, its in-degrees are uniformly bounded, and its out-neighborhoods have polynomial growth, must \(E_{\vec G}\) be hyperfinite?

The growth hypothesis is not fully specified in the supplied record. Interpreting out-neighborhood growth through the forward balls above, possible precise hypotheses include
\[
\forall x\in X\ \exists C_x\geq 1\ \exists k_x\in\Z_{\geq 0}\ \forall n\in\Z_{\geq 0},\qquad |B^+_{\vec G}(x,n)|\leq C_x(n+1)^{k_x},
\]
or the stronger uniform hypothesis
\[
\exists C\geq 1\ \exists k\in\Z_{\geq 0}\ \forall x\in X\ \forall n\in\Z_{\geq 0},\qquad |B^+_{\vec G}(x,n)|\leq C(n+1)^k.
\]
A further possibility requires a common exponent but permits the multiplicative constant to depend on \(x\). The supplied record does not select among these quantifier patterns or explicitly define out-neighborhood growth. These are candidate interpretations requiring review, rather than a choice of a definitive formulation.
\subsection{Short English statement}
The question asks whether a Borel locally finite directed graph with uniformly bounded in-degrees and polynomially growing forward neighborhoods necessarily has a hyperfinite connectedness relation. Hyperfiniteness means that connectedness is an increasing union of Borel equivalence relations with finite classes. The meaning and uniformity of polynomial growth require clarification from the source.
\subsection{Sources}
\begin{itemize}
\item Local UnsolvedMath registry, permanent ID \(801045\), source problem code \texttt{catalogue entry}; title: \emph{Open Problems in Mathematical Logic --- Problem 57}. The registry attributes the record to Barmpalias, Gao et al., with proposed year 2026, and supplies the status \texttt{open}; this status is not independently verified here.
\item Supplied source reference: Barmpalias, Gao et al., \emph{Open Problems in Mathematical Logic} (2026), Problem 57, \url{https://arxiv.org/abs/2608.26628}. The record reports access on 2026-10-16 and extraction by \texttt{arxiv\_latex}. These are supplied provenance details, not a claim of independent source consultation.
\item Supplied upstream revision: \texttt{d371ee917ac9647fbb06cd6fc529deec6f3dda82}.
\end{itemize}
\end{document}
